% Standalone problem document, generated from the adjacent README.md.
% Compile with XeLaTeX. Mathematical content is maintained in Markdown.
\documentclass[11pt,a4paper]{article}
\usepackage[fontsize=10.5pt]{scrextend}
\usepackage[margin=22mm,headheight=15pt,headsep=9mm,footskip=11mm]{geometry}
\usepackage{fontspec}
\setmainfont{texgyrepagella}[Extension=.otf,UprightFont=*-regular,BoldFont=*-bold,ItalicFont=*-italic,BoldItalicFont=*-bolditalic]
\setsansfont{texgyreheros}[Extension=.otf,UprightFont=*-regular,BoldFont=*-bold,ItalicFont=*-italic,BoldItalicFont=*-bolditalic]
\setmonofont{texgyrecursor}[Extension=.otf,UprightFont=*-regular,BoldFont=*-bold,ItalicFont=*-italic,BoldItalicFont=*-bolditalic,Scale=MatchLowercase]
\usepackage{amsmath,amssymb}
\usepackage{unicode-math}
\setmathfont{texgyrepagella-math.otf}
\usepackage{xcolor,microtype,enumitem,needspace,fancyhdr}
\usepackage{bookmark}
\definecolor{ink}{HTML}{17344B}
\definecolor{accent}{HTML}{197D80}
\definecolor{muted}{HTML}{596773}
\hypersetup{colorlinks=true,linkcolor=accent,urlcolor=accent,pdftitle={SP-01 - The
sharp generic threshold for spectral-subspace
rotation},pdfauthor={Open Problems in Numerical Linear Algebra}}
\urlstyle{same}
\setlength{\parindent}{0pt}
\setlength{\parskip}{5pt plus 1pt minus 1pt}
\linespread{1.04}
\setlength{\emergencystretch}{3em}
\widowpenalty=10000
\clubpenalty=10000
\displaywidowpenalty=10000
\allowdisplaybreaks[1]
\setlist{nosep,leftmargin=1.5em}
\providecommand{\tightlist}{\setlength{\itemsep}{0pt}\setlength{\parskip}{0pt}}
\providecommand{\pandocbounded}[1]{#1}
\pagestyle{fancy}
\fancyhf{}
\fancyhead[L]{\sffamily\footnotesize\color{muted}OPEN PROBLEMS IN NUMERICAL LINEAR ALGEBRA}
\fancyhead[R]{\sffamily\bfseries\footnotesize\color{accent}SP-01}
\fancyfoot[L]{\parbox[b]{0.9\textwidth}{\sffamily\fontsize{8}{10}\selectfont\color{muted}Editorial ratings; a bounded search cannot certify that a problem remains unsolved.\\Literature check: 2026-09-08}}
\fancyfoot[R]{\sffamily\footnotesize\color{muted}\thepage}
\renewcommand{\headrulewidth}{0.3pt}
\renewcommand{\headrule}{\hbox to\headwidth{\color{accent}\leaders\hrule height \headrulewidth\hfill}}
\makeatletter
\renewcommand{\section}{\Needspace{5\baselineskip}\@startsection{section}{1}{0pt}{12pt plus 2pt minus 2pt}{4pt}{\sffamily\bfseries\large\color{ink}}}
\renewcommand{\subsection}{\Needspace{4\baselineskip}\@startsection{subsection}{2}{0pt}{9pt plus 1pt minus 1pt}{3pt}{\sffamily\bfseries\normalsize\color{ink}}}
\makeatother
\setcounter{secnumdepth}{0}
\begin{document}
{\sffamily\small\color{muted}Eigenvalues and inverse problems\par}
\vspace{3pt}
{\raggedright\hyphenpenalty=10000\exhyphenpenalty=10000\sffamily\bfseries\fontsize{21}{24}\selectfont\color{ink}The
sharp generic threshold for spectral-subspace rotation\par}
\vspace{7pt}
{\sffamily\small\textbf{Difficulty:} challenging\quad\textbf{Importance:} interesting
to the community\par}
\vspace{4pt}
{\color{accent}\hrule height 0.8pt}
\vspace{6pt}
\textbf{Status:} explicit open constant problem; no later resolution
located

Let \(A\) be a self-adjoint, possibly unbounded operator on a separable
complex Hilbert space. Suppose \(\sigma(A)=\sigma\cup\Sigma\), where the
two nonempty closed sets satisfy
\(d=\operatorname{dist}(\sigma,\Sigma)>0\). For a bounded self-adjoint
perturbation \(V\), set \[
P=E_A(\sigma),\qquad Q=E_{A+V}(O_{d/2}(\sigma)),\qquad
O_r(S)=\{x\in\mathbb R:\operatorname{dist}(x,S)<r\},
\] where \(E_B\) denotes the spectral projection measure of \(B\).

Does \(\|V\|<d/2\) always imply \(\|P-Q\|<1\)? All norms are operator
norms. Equivalently, is the optimal universal constant
\(c_{\mathrm{opt}}=1/2\) for guaranteeing that the maximal angle
\(\arcsin\|P-Q\|\) is strictly below \(\pi/2\)?

The two spectral sets may interlace: no ordering or disjoint-convex-hull
hypothesis is imposed. The source's operator formulation is retained;
Hermitian matrix invariant-subspace perturbation is its
finite-dimensional setting. This problem concerns subspace orientation
after perturbation, beyond preservation of a spectral gap.

\subsection{References}\label{references}

A. Seelmann, \emph{Notes on the subspace perturbation problem for
off-diagonal perturbations}, Proceedings AMS 144 (2016), 3825--3832,
introduction p.3826, the paragraph on general perturbations
(\href{https://arxiv.org/pdf/1412.6294}{primary manuscript};
\href{https://doi.org/10.1090/proc/13118}{journal}). Seelmann, \emph{On
an estimate in the subspace perturbation problem}, Journal d'Analyse
Mathématique 135 (2018), 313--343
(\href{https://arxiv.org/abs/1310.4360}{primary manuscript}). Seelmann,
\emph{Unifying the treatment of indefinite and semidefinite
perturbations in the subspace perturbation problem}, Operators and
Matrices 15 (2021), 1181--1188, Theorems 1.1--1.2
(\href{https://files.ele-math.com/articles/oam-15-74.pdf}{primary
paper}).

\subsection{Status check --- 2026-09-08}\label{status-check-2026-09-08}

The 2021 generic theorem retains the sufficient constant
\(c_{\mathrm{crit}}=0.4548399\ldots\); its sharp theorem at a larger
threshold assumes extra spectral separation. The 2018 paper's solved
optimization problem determines an upper angle bound, not the
conjectured threshold. Searches for ``subspace perturbation'', ``optimal
constant'', ``1/2'', ``Seelmann'', ``conjecture'', and 2024--2026 found
no resolution. No recent explicit reaffirmation of the exact endpoint
conjecture was located; the status evidence is bounded.
\end{document}
